SHYFEM has been coupled to the WAVEWATCH III (WW3) wave model.
The SHYFEM model was coupled to WW3 based on a so called ``hard coupling'', e.g.
binding the models directly without any so called coupling library such as
(ESMF, OASIS, PGMCL or others). The benefit is on the hand, minimal memory
usage, very limited code changes in both models. Top-level approach using the
coupled model framework based on having a routine for initialization,
computation and finalization, that allows a neat inclusion of WW3 in any kind
of flow model.

\paragraph{Installation and compilation}

The coupled model is distributed within the SHYFEM distribution, however the WW3 version
compliant to SHYFEM must by asked at shyfem@gmail.com.

In order to compile the coupled SHYFEM-WW3 model, you have to use Intel fotran compiler:
\begin{code}
FORTRAN_COMPILER = INTEL
C_COMPILER = INTEL
\end{code}
The variable |WW3| in the file |Rules.make| should be set as:
\begin{code}
WW3 = true
\end{code}
Additionally the |WW3DIR| should point to the folder containing the WW3 code, e.g.:
\begin{code}
WW3DIR = \$(HOME)/path/to/WW3). 
\end{code}
SHYFEM should be compiled for running in parallel by setting 
|PARALLEL_MPI = NODE|.

Moreover, WW3 also needs additional libraries:
\begin{itemize}
\item |netcdf| compiled with the same compiler.  You must set |NETCDF = true|
and specify the variable |NETCDFDIR| to indicate the directory where the 
libraries and its include files can be found.
\item METIS and PARMETIS compiled with the same compiler. You must set
|PARTS = PARMETIS| and set the variables |METISDIR| and |PARMETISDIR|
to indicate the directory where the libraries and its include files can be
found.
\end{itemize}
At the end the Rules.make should be modified as follows:
\begin{code}
PARALLEL_MPI = NODE
PARTS = PARMETIS
METISDIR = /path/to/metis
PARMETISDIR = /path/to/parmetis
NETCDF = TRUE
NETCDFDIR = /path/to/netcdf
\end{code}

The compilation of WW3 can be customized changing a set of model options,
defined with the variable |WW3CFLAGS| in the file |fem3d/Makefile|. They
correspond to the parameters listed in the file |switch| needed by the stand 
alone WW3 installation. For a detailed description of the options see the WW3 
manual (section 5.9).

You can now follow the general SHYFEM installation instruction to compile the 
code:
\begin{code}
make cleanall
make
\end{code}
To run WW3 it is however required to compile the stand alone WW3 
model. This step is needed for having WW3 source code and pre- and post-processing 
tools (e.g., ww3\_grid). 
It is important that the so called ``switch'' file of WW3 (see WW3 documentation) 
contains the same parameters listed with variable |WW3CFLAGS| are identical 
for both modules. The default setup provided with |WW3CFLAGS| is basically 
identical to the setup of SHOM and Meteo France as well as the USACE based 
on implicit time stepping on unstructured grids. Basically, there is no need 
for modification of the ``switch'' file with the given settings but all 
settings are supported as long the unstructured grid option in WW3 is used.

For installing and comping the stand alone WW3 model you can follow
the following procedure (see the WW3 manual for more informations):
\begin{enumerate}
%\item download the ww3\_shyfem branch from the ERDC git repository: 
%git clone --branch ww3\_shyfem https://github.com/erdc/WW3/tree/ww3\_shyfem;
\item move to the WW3 directory;
\item set the NetCDF path with:
\begin{code}
export WWATCH3_NETCDF=NC4
export NETCDF_CONFIG=/netcdf-dir/nc-config
\end{code}
\item setup the model using
\begin{code}
./model/bin/w3_setup /home/model/bin/WW3/model -c <cmplr> -s <swtch>
\end{code}
 with |cmplr| the compiler |comp| and |link| files (e.g., shyfem
for files comp.shyfem and link.shyfem) and |swtch| the switch file (e.g., shyfem 
for a switch file located in ./model/bin having name switch\_shyfem).
\item compile WW3 with:
\begin{code}
./w3_automake ww3_grid
\end{code}
the compiled programms are located in |./model/exe/ww3_grid|
\end{enumerate}

\paragraph{Coupling description}
The implementation was done by developing two modules, one in
SHYFEM (subww3.f) and the other one in WW3 (w3cplshymfen.F90). The 
SHYFEM coupling module is connected by the ``use'' statement to the 
modules of the WW3 code. In this way one can access all fields from 
the WW3 in SHYFEM and vice versa. In the WW3 module all spectral based
quantities are computed and SHYFEM is directly assigning the various 
arrays, based on global arrays, to the certain domains for each of the 
model decompositions.

In terms of WW3 we have used the highest-level implementation based
on the WW3 multigrid driver. This allows basically to couple any kind
of grid type with SHYFEM (rectangular, curvi-linear, SMC and
unstructured). It offers of course the possibility to run WW3 based
on a multigrid SETUP and coupled to SHYFEM. Both models use their
native input files for running the model except that the forcing by
wind within the coupled model is done via SHYFEM and the flow field
is by definition provided based on the coupling to SHYFEM. In this
way both models are fully compatible to the available documentation.

The two numerical models (SHYFEM and WW3) should exchange all the
variables that are needed to simulate the current-wave interaction.
The following variables sare passed by SHYFEM to WW3:
\begin{itemize}
\item water levels;
\item three-dimensional water currents;
\item wind components.
\end{itemize}

The following physical quantities have been computed in WW3 and are 
available for SHYFEM:
\begin{itemize}
\item significant wave height;
\item mean wave period;
\item peak wave period;
\item main wave direction (the where the wave go);
\item radiation stress,
\item wave pressure;
\item wind drag coefficient;
\item Stokes velocities.
\end{itemize}

The exchange module of SHYFEM contains all the needed infrastructure
for initializing, calling and finalizing the wave model run. Two
subroutines (getvarSHYFEM and getvarWW3), which are called before and
after the call to the wave model in order to obtain currents, water
levels and on the other hand integral wave parameters and radiation
stresses.

The significant wave height, mean wave period and main wave 
direction are written by SHYFEM in the |.wave.shy| file according
to the values of the variables |idtwav| and |itmwav| in the |wave|
section of the SHYFEM parameter file (see below).

\paragraph{Running the coupled SHYFEM-WW3 model}
Since the wave output are written by SHYFEM, the output are
set to off in WW3, which reduces disk usage and significantly reduces 
the parallel overhead due to the output part. 

As we have utilized here the implicit scheme in WW3, which was
developed by Roland \& Partner and the implicit scheme is well
validated and unconditionally stable. The choice of the time step
should be according to the physical time scales of the modelled
processes. 

In order to run the coupled model, we suggest the following
procedure:
\begin{enumerate}
\item Setup the SHYFEM set-up for the region of interest. 
SHYFEM needs a mesh in the .grd format. Preprocess the SHYFEM grid with a selected number of domains 
without bandwidth optimization with
\begin{code}
shypre -noopti grid.grd|
\end{code}
SHYFEM need a parameter file and all forcing files (see SHYFEM documentation). In the SHYFEM 
parameter file the section |\$waves| must be activated:
\begin{code}
        iwave  = 11             dtwave = '5m'
        itmwav = '2012-02-01::00:00'    idtwav = '30m'
\end{code}
The time step for coupling with WW3 |dtwave| must be set equal
to the SHYFEM and WW3 timesteps.
|idtwav|, |itmwav| should also be set for determining the 
time step and start time for writing to the output file |.wave.shy| .

\item Setup the WW3 model for the region of interest. In the
running directoty, WW3 requires:
\begin{enumerate}
\item the numerical mesh in the GMSH format .msh. The SHYFEM
mesh in .grd format can be easily converted to .msh with the command:
\begin{code}
shybas -msh grid.bas
\end{code}
it creates a bas.msh file ready to be used by WW3.
\item the file |ww3_grid.nml| containing the spectrum, run, timesteps
and grid parameterizations (see WW3 documentation). This file also
contains the name of mesh file previously converted in .msh format and
namelist files.
\item the namelist file containing the variables for setting
the tunable parameters for source terms, propagation schemes, and 
numerics.
\item the file |ww3_multi.nml| which handles the time steps and 
the I/O. Since WW3 receives flow and wind from SHYFEM there are 
no input fields that need to be prescribed. The output part of the 
wave model itself is handled by SHYFEM as well.
\end{enumerate}

Before running the model, the setup |ww3_grid| tool (found in |WW3/build/bin|
directory) needs to be run. The output, 
which is named per default ``mod\_def.ww3'' needs to be renamed 
with the extension defined in parameter
\begin{code}
MODEL(1)\%NAME = 'med'
\end{code}
in |ww3_multi.nml| (e.g., mod\_def.med in the above mentioned case). 
The |ww3_grid| tool must be run every time the namelist or grid files 
are modified. 
\begin{code}
./ww3_grid
\end{code}

\item run the coupled model with the |shyfem| command. Since the code
is optimized to run with MPI on domain decomposition, we suggest to
run the coupled model in parallel using
\begin{code}
mprun -np np shyfem namelist.str
\end{code}
with np the number of processors.

\end{enumerate}

\subsection{WW3 Configuration for New SHYFEM Architecture}

This section documents the WW3 coupling configuration for the current SHYFEM
architecture (version 8.1.1 and later). The coupling implementation has been
ported from the previous SHYFEM version with modifications to accommodate
architectural changes in the core SHYFEM codebase.

\paragraph{Overview of Changes}

The new SHYFEM architecture includes the following key changes affecting WW3 coupling:
\begin{itemize}
\item The |WW3SWITCH| variable allows dynamic selection of WW3 configuration files
\item The |shympi_aux| module has been merged into the |shympi| module
\item WW3 switch files are now loaded dynamically from the WW3 repository
\item BLAS function conflicts are prevented using W3\_SHYFEM preprocessor guards
\end{itemize}

\paragraph{Prerequisites}

Before configuring SHYFEM with WW3 coupling, ensure the following dependencies
are installed:
\begin{itemize}
\item MPI library (OpenMPI or MPICH)
\item METIS library (for graph partitioning)
\item PARMETIS library (for parallel graph partitioning)
\item NetCDF library with Fortran bindings (must be compiled with the same
Fortran compiler as SHYFEM)
\item Git (for cloning the WW3 repository)
\end{itemize}

\paragraph{WW3 Source Installation}

Users must obtain and configure the WW3 source code separately:

\subparagraph{Step 1: Clone WW3 from GitHub}
\begin{code}
git clone https://github.com/erdc/WW3.git
cd WW3
\end{code}

The ERDC WW3 repository is the official source for WAVEWATCH III.

\subparagraph{Step 2: Create or Modify WW3 Switch File}

Edit the switch file that will be used for compilation (e.g.,
|model/bin/switch_itedev|). The switch file must include the |SHYFEM| flag
to enable SHYFEM coupling:

\begin{code}
NOGRB PDLIB METIS SCOTCH PR3 UQ FLX0 LN1 ST4 STAB3 NL1 BT1 DB1 MLIM TR0 BS0 IC0 IS0 REF0 XX0 WNT2 WNX1 CRT1 CRX1 O0 O1 O2 O2a O2b O2c O3 O4 O5 O6 O7 SHYFEM
\end{code}

The |SHYFEM| flag at the end is critical for proper coupling.

\subparagraph{Step 3: Prevent BLAS Function Conflicts}

The WW3 file |w3profsmd.F90| contains BLAS functions (DNRM2, DDOT, DAXPY) that
conflict with SHYFEM's libsparskit.a library. These functions must be wrapped
with preprocessor guards:

Edit |model/src/w3profsmd.F90| and add the following:

Around line 4498, add:
\begin{code}
#ifndef W3_SHYFEM
\end{code}

Around line 4688 (after the DAXPY subroutine), add:
\begin{code}
#endif
\end{code}

This wraps the BLAS functions DNRM2, DDOT, DAXPY, and DLASSQ so they are
excluded when compiling with the W3\_SHYFEM flag.

\subparagraph{Step 4: Obtain SHYFEM Coupling Interface}

The file |w3cplshyfem.F90| contains the coupling interface between SHYFEM and
WW3. This file is NOT included in the standard WW3 distribution and must be
obtained separately:

\begin{itemize}
\item Contact: shyfem@gmail.com
\item Or from the old SHYFEM distribution at: |/path/to/old/shyfem/WW3/model/src/w3cplshyfem.F90|
\end{itemize}

Copy this file to: |$WW3DIR/model/src/w3cplshyfem.F90|

\paragraph{SHYFEM Configuration}

Configure the SHYFEM |Rules.make| file with the following settings:

\subparagraph{Parallel MPI Configuration}
\begin{code}
PARALLEL_MPI = NODE
\end{code}
WW3 coupling requires MPI parallelization to be enabled.

\subparagraph{Partitioning Configuration}
\begin{code}
PARTS = PARMETIS
METISDIR = /path/to/metis
PARMETISDIR = /path/to/parmetis
\end{code}
Replace |/path/to/metis| and |/path/to/parmetis| with the actual installation
paths. These can often be found using:
\begin{code}
locate libmetis.a
locate libparmetis.a
\end{code}

\subparagraph{NetCDF Configuration}
\begin{code}
NETCDF = true
NETCDFDIR = /path/to/netcdf
NETCDFFDIR = /path/to/netcdf-fortran
\end{code}
The NetCDF paths can be found using:
\begin{code}
nc-config --prefix
nf-config --prefix
\end{code}
Ensure NetCDF was compiled with the same Fortran compiler as SHYFEM.

\subparagraph{WW3 Configuration}
\begin{code}
WW3 = true
WW3DIR = /path/to/WW3
WW3SWITCH = switch_itedev
\end{code}

Key points:
\begin{itemize}
\item |WW3DIR| should point to the base directory of the cloned WW3 repository
\item |WW3SWITCH| specifies the switch file name in |$WW3DIR/model/bin/|
\item The switch file must include the |SHYFEM| flag (see Step 2 above)
\item This variable is new in the current SHYFEM architecture
\end{itemize}

\paragraph{Complete Rules.make Example}

A complete configuration example:
\begin{code}
FORTRAN_COMPILER = GFORTRAN    # or INTEL
C_COMPILER = GNU               # or INTEL

PARALLEL_MPI = NODE
PARTS = PARMETIS
METISDIR = /usr/local/lib
PARMETISDIR = /usr/local/lib

NETCDF = true
NETCDFDIR = /home/user/opt/netcdf_gfortran
NETCDFFDIR = /home/user/opt/netcdf_gfortran

WW3 = true
WW3DIR = /home/user/git/WW3
WW3SWITCH = switch_itedev
\end{code}

\paragraph{Compilation}

After configuration, compile SHYFEM:
\begin{code}
make cleanall
make
\end{code}

The build system will:
\begin{enumerate}
\item Read the switch file from |$WW3DIR/model/bin/$WW3SWITCH|
\item Extract WW3 switches and define them as preprocessor flags
\item Compile WW3 coupling code with W3\_SHYFEM flag
\item Link against WW3, PARMETIS, METIS, and NetCDF libraries
\end{enumerate}

Verify the executable was created:
\begin{code}
ls -lh bin/shyfem
\end{code}
The executable should be approximately 14 MB in size.

\paragraph{Key Differences from Previous SHYFEM}

For users migrating from older SHYFEM versions:

\begin{itemize}
\item \textbf{WW3SWITCH variable}: New variable in |Rules.make| for selecting
WW3 configuration files dynamically
\item \textbf{shympi\_aux module}: This module no longer exists; functionality
has been merged into the |shympi| module
\item \textbf{WW3 source modifications required}: Users must modify the WW3
source code (switch file and w3profsmd.F90) before compilation
\item \textbf{Dynamic switch loading}: WW3 switches are loaded from the
repository at compile time rather than hardcoded
\item \textbf{No WW3 distribution}: WW3 source is no longer distributed with
SHYFEM; users must clone from GitHub
\end{itemize}

\paragraph{Verification}

A test case |setup_ww3| is available for verifying the WW3 coupling:
\begin{code}
cd tests/setup_ww3
# Follow test case instructions
\end{code}

This test case provides a complete example of SHYFEM-WW3 coupling configuration
and can be used to verify that the compilation and coupling are working correctly.

\paragraph{Troubleshooting}

Common issues and solutions:

\begin{itemize}
\item \textbf{BLAS symbol conflicts}: Ensure |w3profsmd.F90| has been modified
with |#ifndef W3_SHYFEM| guards
\item \textbf{Missing w3cplshyfem.F90}: Contact shyfem@gmail.com to obtain
this file
\item \textbf{Switch file not found}: Verify |WW3SWITCH| matches the filename
in |$WW3DIR/model/bin/|
\item \textbf{NetCDF errors}: Ensure NetCDF was compiled with the same Fortran
compiler
\item \textbf{METIS/PARMETIS not found}: Check |METISDIR| and |PARMETISDIR|
point to directories containing |lib/| subdirectories
\end{itemize}

For additional support, contact: shyfem@gmail.com
